\documentclass[11pt,a4paper]{article}
\usepackage[T1]{fontenc}
\usepackage[utf8]{inputenc}
\usepackage[spanish]{babel}
\usepackage{lmodern}
\usepackage[a4paper,left=2cm,right=2cm,top=2.3cm,bottom=2.2cm]{geometry}
\usepackage{tabularx,array,hyperref}
\hypersetup{colorlinks=true,urlcolor=blue,pdftitle={Portada SGA - EV-PRAC-01},pdfauthor={Herrera Moran Gilmar Javier}}
\setlength{\parindent}{0pt}
\renewcommand{\arraystretch}{1.55}
\pagestyle{empty}
\begin{document}
\begin{center}
{\large UNIVERSIDAD T\'ECNICA ESTATAL DE QUEVEDO}\\[6pt]
{\normalsize Facultad de Ciencias de la Computaci\'on y Dise\~no Digital}\\[3pt]
{\normalsize Carrera de Ingenier\'ia de Software}\\[14pt]
{\normalsize Asignatura: Ingenier\'ia de Requerimientos (ISR-401)}\\[42pt]
{\LARGE\bfseries Evaluaci\'on de la Unidad I\\[4pt] Fundamentos de Ingenier\'ia de Requisitos}\\[14pt]
{\large\bfseries EV-PRAC-01}\\[14pt]
{\normalsize Prueba pr\'actica escrita a mano y digitalizaci\'on en \LaTeX}\\[8pt]
{\normalsize Caso de estudio: Sistema de Monitoreo de la Cadena de Fr\'io <<FR\'IOVAC>>}
\end{center}
\vspace{1.4cm}
\begin{tabularx}{\textwidth}{|>{\bfseries}p{4.5cm}|X|}\hline
Estudiante: & Herrera Moran Gilmar Javier \\ \hline
Nivel y paralelo: & Cuarto nivel, Software A \\ \hline
Docente: & PhD. Gleiston Cicer\'on Guerrero Ulloa \\ \hline
C\'odigo de la tarea: & EV-PRAC-01 \\ \hline
Fecha: & Mi\'ercoles 18 de noviembre de 2026 \\ \hline
\end{tabularx}
\vspace{1.3cm}
\begin{center}
{\bfseries Repositorio GitHub}\\[12pt]
% URL completa en una sola l\'inea, sin saltos de l\'inea.
\mbox{\href{https://github.com/GilmaR2006/ISR401-HerreraMoranGilmar}{\small\texttt{https://github.com/GilmaR2006/ISR401-HerreraMoranGilmar}}}
\end{center}
\vfill
\begin{center}
Quevedo, Ecuador --- 2026
\end{center}
\end{document}
